\section{First inverse decomposition} \label{Section: descomposición inversa}



Let $\lambda\in X^+$. Our goal in this section is to express $\bfM_{\lambda}^{\succeq \alpha_{i,j}}$ as a linear combination of terms of the form $\bfM_{\mu}^{\succ \alpha_{i,j}}$. 
Since $\Phi^{\succeq \alpha_{i,j}}=\Phi^{\succ \alpha_{i,j}}\cup\{\alpha_{i,j}\}$, we have
\begin{equation}
   \bfM_{\lambda}^{\succeq \alpha_{i,j}}
   \;=\;
   \bfM_{\lambda}^{\succ \alpha_{i,j}}
   \;-\;
   q\,\bfM_{\lambda-\alpha_{i,j}}^{\succ \alpha_{i,j}}.
\end{equation}
If $\lambda-\alpha_{i,j}\in X^+$ then this equality already gives the desired decomposition. 
Otherwise, a refinement is required for proving the positivity conjecture.
In \Cref{prop: second version} we provide an initial decomposition, and a sharper form is established in \Cref{sec: second inverse decomp} (see \Cref{prop: second version refinado}). 
The latter is the version best suited for proving the positivity conjecture.




\subsection{$P$ and $Q$ operators}
In this section, we introduce the notation needed to express the first of the aforementioned decompositions. The main tools are the $P$ and $Q$ operators, which act on the set of weights $X$.



\begin{definition}\rm 
Let $\lambda\in X$ and $h $ be an  integer greater than or equal to $2$. 
We say that 
\begin{enumerate}[(a)]
    \item  $\lambda$ is $L$-dominant at position $r$ and height $h$ if $\langle \lambda , \alpha_k \rangle \geq 0$ for all  $1\leq k\leq r$ with $k\equiv r \mod{h}$. 
    \item  $\lambda$ is $R$-dominant at position $r$ and height $h$ if $\langle \lambda , \alpha_k \rangle \geq 0$ for all $r \leq k \leq n$ with $k\equiv r \mod{h-1}$.
    \item $\lambda$ is $L$-dominant at position $r$  if $\langle \lambda , \alpha_k \rangle \geq 0$ for all  $1\leq k\leq r$. 
    \item $\lambda$ is $R$-dominant at position $r$ if $\langle \lambda , \alpha_k \rangle \geq 0$ for all $r \leq k \leq n$.
\end{enumerate}
We denote by $X_{r,h}^+(L) $ (resp. $X_{r,h}^+(R)$) the set of all weights $L$-dominant (resp. $R$-dominant) at position $r$ and height $h$.
Similarly, we denote by $X_{r}^+(L) $ (resp. $X_{r}^+(R)$) the set of all weights $L$-dominant (resp. $R$-dominant) at position $r$.
\end{definition}



\begin{definition}\rm \label{def. p and q}
Let $\lambda \in X$, $1\leq i<j \leq n$ and $h=j-i+1$. 
We define
\begin{equation}
\begin{array}{ll}
   \varrho = \varrho_{j,h}(\lambda )  = &   \displaystyle  \max \{   1\leq k \leq i \mid k\equiv i \mod h \quad \mbox{ and } \quad \lambda -\alpha_{k,j} \in X_{i,h}^+(L)     \}, \\[3pt]
 \vartheta=   \vartheta_{j,h}(\lambda)  = &   \displaystyle \min \{ j< k \leq n \mid k\equiv j \mod (h -1) \quad \mbox{ and } \quad  \lambda -\alpha_{j+1,k } \in X_{j,h}^+(R)     \}.
\end{array}
\end{equation} 


\begin{enumerate}[(a)]
    \item If $\varrho$ exists we define $\OP{j,h}(\lambda)= \lambda-\alpha_{\varrho,j}$. Otherwise, we say $\OP{j,h}(\lambda)$ is not defined. 
    
    We call $\varrho$ the integer associated to $\OP{j,h}(\lambda)$.
    \item If $\vartheta$ exists we define $\OQ{j,h}(\lambda) =\lambda-\alpha_{j+1,\vartheta}$. Otherwise, we say $\OQ{j,h}(\lambda)$ is not defined. 
    
    We call $\vartheta $ the integer associated to $\OQ{j,h}(\lambda)$.
\end{enumerate}
In order to simplify notation, if $h$ is clear for the context, we just write
$ \OP{j}(\lambda) =\OP{j,h}(\lambda)  $ and $\OQ{j}(\lambda)=\OQ{j,h}(\lambda) $.
\end{definition}


\begin{example}\rm
Let $n=9$, $\lambda=[1,1,0,0,0,0,0,1,1]$, $j=7$ and  $h=3$. 
Then, $\varrho =2$, $\vartheta = 9$, 
    $P_{7,3}(\lambda) = \lambda -\alpha_{2,7} = [2,0,0,0,0,0,-1,2,1] $ and  $Q_{7,3}(\lambda) = \lambda - \alpha_{8,9}= [1,1,0,0,0,0,1,0,0]$.
\end{example}


\begin{definition}\rm\label{def: words over weights} 
    Let $\lambda \in X$, $1\leq j \leq n $ and $h\geq 2$. 
    For $k\geq 1$ we define
    \begin{equation}
    \begin{array}{rl} 
        P^k_{j,h} (\lambda) = &\OP{j-k+1,h}\OP{j-k+2,h}\cdots \OP{j-1,h}\OP{j,h}(\lambda).\\ [3pt]
        Q^{k}_{j,h} (\lambda) =& \OQ{j,h}\OQ{j-1,h}\cdots \OQ{j-k+1,h}(\lambda).\\[3pt]
        Q^{-k}_{j,h} (\lambda) =& \OQ{j+k-1,h}\OQ{j+k-2,h}\cdots \OQ{j+1,h}\OQ{j,h}(\lambda).\\[3pt]
        v^k_{j,h}(\lambda) =&  Q^{k}_{j,h} P^k_{j,h} (\lambda).
        % (v_{k}P^{m})_{j,h}(\lambda) =& (v_{k})_{(j-m,h)}(\OP{}^{m})_{(j,h)}(\lambda).\\[2pt]
        % (P^{m}v_{k})_{j,h}(\lambda) =& (\OP{}^{m})_{(j-k,h)}(v_{k})_{(j,h)}(\lambda).
    \end{array}
    \end{equation}
  By convention, we define all these operators for $k=0$ simply as $\lambda$. 
\end{definition}
\begin{remark}\rm
    Observe that in \Cref{def: words over weights}, the definitions of the operators $Q^{k}_{j,h}$ and $Q^{-k}_{j,h}$ are equivalent up to a shift in the index. More precisely, one has
    \begin{equation}
        Q^{k}_{j,h}(\lambda) = Q^{-k}_{j-k+1,h}(\lambda).
    \end{equation}
    Throughout the text, we will freely use both notations depending on convenience of reading or context.
\end{remark}

    
We stress that the operators defined before correspond to a sequence of compositions of $P$ and $Q$ operators. Throughout these sequence it may well be the case that one of the elements involved in the composition is not defined. In this case, we simply say that the whole operator is not defined in a given weight. 

\begin{example}\rm
Let $(i,j)=(5,8)$ and $h = 4$. Consider the weight $\lambda = [2,0,1,2,0,0,0,0,0,0,1,0,1,0]_{\varpi}$. 
Applying our operators, we obtain:
\begin{equation}
\begin{aligned}
    P_{8,4}^2(\lambda) &= [1, 0, 2, 1, 0, 0, -1, 0, 1, 0, 1, 0, 1, 0]_{\varpi}, \\[5pt]
    v_{8,4}^2(\lambda) = Q_{8,4}^2P_{8,4}^2(\lambda) &= [1, 0, 2, 1, 0, 0, 0, 0, 0, 0, 0, 1, 0, 1]_{\varpi}, \\[5pt]
    P_{6,4}^2v_{8,4}^2(\lambda) &= [2, 0, 1, 1, -1, 0, 1, 0, 0, 0, 0, 1, 0, 1]_{\varpi}.
\end{aligned}
\end{equation}
This allows us to conclude that the operators $P_{8,4}^2$, $v_{8,4}^2$, and $P_{6,4}^2v_{8,4}^2$ are all defined at $\lambda$.

\medskip
On the other hand, if we change the initial weight to $\lambda = [2,0,1,2,2,0,0,0,0,0,0,0,1,0]_{\varpi}$, we compute:
\begin{equation}
    P_{8,4}^2(\lambda) = [2, 0, 2, 2, 1, 0, -1, 0, 1, 0, 0, 0, 1, 0]_{\varpi}.
\end{equation}
Thus, $P_{8,4}^2(\lambda)$ is defined. However, notice that $\pint{P_{8,4}^2(\lambda)}{\alpha_{11}} = \pint{P_{8,4}^2(\lambda)}{\alpha_{14}} = 0$. This implies that $Q_{7,4}(P_{8,4}^2(\lambda))$ is undefined, and consequently, neither $v_{8,4}^2(\lambda)$ nor $P_{6,4}^2v_{8,4}^2(\lambda)$ are defined.
\end{example}




The following lemma shows certain relationship between the integers associated to the different $P$ operators that occur in the composed operator $P^k$. 


\begin{lemma} \label{lem: desigualdad de varrhos}
    Let $\lambda \in X^+_{j}(L)$, $\alpha_{i,j}\in \Phi^{\geq 2}$ and $h=j-i+1$. 
    For $r\geq 0$ we set $\lambda_{r} = P^r_{j,h}   (\lambda) $ and  $\varrho_{r}=\varrho_{j-r+1,h}(\lambda_{r-1})   $.
    Suppose that  $\lambda_{k}$ is defined for some  $k\geq 1$ and that  $\varrho_{k}\geq 2$. 
    Then, $\lambda_{k+1}$ is defined and $\varrho_{k+1} \geq \varrho_{k} - 1$.
    In particular, if $\lambda_{k}$ is defined but $\lambda_{k+1}$ is not then $\varrho_k=1$.   
\end{lemma}

\begin{proof}
By definition $\varrho_k\equiv i-k +1  \mod h $. 
It follows that $\varrho_k -1\equiv i-k   \mod h $.
Since $\lambda \in X_j^+(L) $  and $\lambda_{k-1} $ is defined we have $\lambda_{k-1} \in X^{+}_{j-k+1}(L)$. 
In particular, $\pint{\lambda_{k-1}}{\alpha_{\varrho_k-1}}\geq 0$. It follows that
\begin{equation}
    \pint{\lambda_{k}}{\alpha_{\varrho_k-1}} = \pint{\lambda_{k-1}-\alpha_{\varrho_k,j-k+1}}{\alpha_{\varrho_k-1}} = \pint{\lambda_{k-1}}{\alpha_{\varrho_k-1}} -\pint{\alpha_{\varrho_k,j-k+1}}{\alpha_{\varrho_k-1}} = \pint{\lambda_{k-1}}{\alpha_{\varrho_k-1}} +1 \geq 1. 
\end{equation}
Thus  $\lambda_k - \alpha_{\varrho_k-1, j-k+1}\in X^+_{j-k,h}(L) $. 
We conclude that  $\lambda_{k+1}$ is defined and $\varrho_{k+1} \geq \varrho_{k} - 1$.
\end{proof}


\begin{corollary} \label{lem: ceros por Ps}
 Let $\lambda \in X^+_{j}(L)$ and $\alpha_{i,j}\in \Phi^{\geq 2}$ and $h=j-i+1$.
 For $r\geq 0$ we set $\lambda_{r} = P^r_{j,h}   (\lambda) $ and  $\varrho_{r}=\varrho_{j-r+1,h}(\lambda_{r-1})  $.
    Suppose that  $\lambda_{k}$ is defined for some  $1\leq k \leq h$. 
    Then, we have
    $\pint{\lambda_{k-1}}{\alpha_{t}} = 0$,
    for all $\displaystyle  t\in U_d\coloneqq \bigcup_{m=0}^{d} [i-k+1-mh,i-mh]$, where $d$ is the unique integer such that $\varrho_k = i-k+1 -(d+1)h$.
\end{corollary}

\begin{proof}
If $d=-1$ then $U_d=\emptyset$ and the statement is empty, so assume $d\geq 0$.

We argue by contradiction. Suppose that there exists $t\in U_d$ with 
$\pint{\lambda_{k-1}}{\alpha_t}\neq 0$.
Since $U_d\subseteq [\varrho_k+h,i]$, we have
\begin{equation} \label{eq: t is in  between}
    \varrho_k+h \le t \le i.
\end{equation}

Because $\lambda\in X_j^+(L)$, we have $\lambda_{k-1} \in X_{j-k+1}^+(L)$.  
Moreover, since $k\le h$, it follows that $t\le i\le j-k+1$, and hence
$\pint{\lambda_{k-1}}{\alpha_t}>0$.
The inequality $\varrho_k+h \le t$ and the maximality of $\varrho_k$ imply that 
$t\not\equiv i-k+1 \pmod{h}$.  
Since $t\in U_d$, there exists $0\le s<k-1$ such that 
$t\equiv i-s \pmod{h}$.
By maximality of $\varrho_{s+1}$ we have $t\le \varrho_{s+1}$.
Applying \Cref{lem: desigualdad de varrhos} repeatedly gives
\begin{equation}
    \varrho_{s+1} \le \varrho_k + (k-s-1) < \varrho_k + h.
\end{equation}
Therefore,
    $t < \varrho_k + h$,
contradicting \eqref{eq: t is in  between}.
This contradiction proves the lemma.
\end{proof}

  

Both \Cref{lem: desigualdad de varrhos} and \Cref{lem: ceros por Ps} have their $Q$-counterpart that we now enunciate. 
In both cases we omit the proof as it is almost identical to the one of its $P$-counterpart. 


\begin{lemma} \label{lem: desigualdad de varthetas}
    Let $\mu \in X^+$ and $\alpha_{i,j}\in \Phi^{\geq 2}$ and $h=j-i+1$. 
    For $r\geq 0$ we set $\mu_{r} = Q^{r}_{i,h}(\mu)$ and $\vartheta_{r}=\vartheta_{i+r,h}(\mu_{r-1})$.
    % Let $\vartheta_{r}$ be the integer associated to $Q_{i+r,h}(\mu_{r-1})$.
    Suppose that $\mu_{k}$ is defined for some integer $k\geq 1$ and $\vartheta_{k} \leq n-1$.
    Then, $\mu_{k+1}$ is defined and $\vartheta_{k+1} \leq \vartheta_{k} + 1$.
% \begin{equation}\label{eq: desigualdad de los varthetas}
%     \vartheta_{k+1} \leq \vartheta_{k} + 1.
% \end{equation}
In particular, if $\mu_{k}$ is defined but $\mu_{k+1}$ is not then $\vartheta_k =n$. 
\end{lemma}




\begin{corollary} \label{lem: ceros por Qs}  
% \David{Revisar si el caso $k=h$ funciona aqu\'i y en caso de ser as\'i, incluirlo aunque no se use mas adelante (por completitud)}
    Let $\alpha_{i,j}\in \Phi^{\geq 2}$ and $h=j-i+1$. 
    Let $r$ be an integer such that $i+1\leq r\leq j$.
    Let $\mu \in X^+_{r+1}(R)$.
    For $x\geq 0$, we set $\mu_{x} = Q^{-x}_{r,h}(\mu)$.
    % Set $\mu_{0} = \lambda$ and define $\mu_{r}\coloneqq Q_{p+r-1,h}(\mu_{r-1})$ for $r\geq 1$.
    Suppose $\mu_k$ is defined for some $1\leq k\leq h$.
    Then, we have $\pint{\mu_{k-1}}{\alpha_{t}} = 0$, for all $t\in U_d$, where
    \begin{equation}\label{eq: Ud de zeros en Q}
     U_{d} =   \bigcup_{m=1}^{d} [r + m(h-1), r + k -1 + m(h-1)]
    \end{equation}
  and  $d$ is the unique integer satisfying $\vartheta_{r+k-1,h}(\mu_{k-1}) = r+k-1 + (d+1)(h-1)$.  
\end{corollary}



The following lemma provides the integers associated to the $Q$ operators involved in the  operator $v^{h-1}_{j,h}$.



\begin{lemma}\label{lem: thetas para el caso h-1}
    Let $\lambda \in X^+$ and $\alpha_{i,j}\in \Phi^{\geq 2}$ and $h=j-i+1$. 
    Suppose $v^{h-1}_{j,h}(\lambda)$ is defined.
    We define
    \begin{equation}\label{eq: varthetas en h-1}
        \vartheta_{r}= \left\{  \begin{array}{ll}
          \vartheta_{i+1,h}(P^{h-1}_{j,h}(\lambda)),   & \mbox{if } r=1;  \\[3pt]
        \vartheta_{i+r,h}(Q^{r-1}_{i+r-1,h}(P^{h-1}_{j,h}(\lambda))),     &  \mbox{if } 2\leq r \leq h-1.
        \end{array}  \right.
    \end{equation}
    Then,  $\vartheta_{r}=j+r$, for $1\leq r\leq h-1 $.
\end{lemma}


\begin{proof}
    We begin by noticing that 
    \begin{equation}
        P^{h-1}_{j,h}(\lambda) = \lambda - \sum_{r=1}^{h-1}\alpha_{\varrho_{r},j-r+1},
    \end{equation}
    where $\varrho_{r} = \varrho_{j-r+1,h}(P^{r-1}_{j,h}(\lambda))$,  for $1\leq r\leq h-1$.
    Using this, we can compute the following: 
    \begin{equation}\label{eq: calculo ph-1}
        \pint{P^{h-1}_{j,h}(\lambda)}{\alpha_{j+1}} = \pint{\lambda}{\alpha_{j+1}} - \sum_{r=1}^{h-1} \pint{\alpha_{\varrho_{r},j-r+1}}{\alpha_{j+1}} = \pint{\lambda}{\alpha_{j+1}} + 1 \geq 1,
    \end{equation}
    where the last inequality follows since $\lambda \in X^+$.


   We recall that $\vartheta_1$ is, by definition, the minimal integer satisfying $P^{h-1}_{j,h}(\lambda)-\alpha_{i+1,\vartheta_1} \in X_{i+1,h}^+(R)$ and $\vartheta_1\equiv i+1 \mod h-1$.
   As $j+1 = i+1 + h-1$, \eqref{eq: calculo ph-1} implies that   $\vartheta_{1} = j+1$, which is our claim for the case $r=1$. 
   

On the other hand,  \Cref{lem: desigualdad de varthetas} yields $\vartheta_{r} \leq \vartheta_{r-1}+1$ for all $2\leq r\leq h-1$.
Furthermore, by \cref{eq: varthetas en h-1},    $ i+r+h-1  \leq \vartheta_r$. 
By combining these two inequalities with the definition of $h$, we obtain 
\begin{equation}
    j+r\leq \vartheta_r\leq \vartheta_{r-1}+1. 
\end{equation}
Finally, an inductive argument with the equality $\vartheta_1=j+1$ yields $\vartheta_r=j+r$, as we wanted to show. 
\end{proof}





We now define the \textbf{degree} of the action of a $P$ or $Q$ operator on a weight $\lambda$. 




\begin{definition}\rm  \label{def: degree of P and Q}
Let $\lambda \in X$, $1\leq j \leq n$ and $h \geq 2$. 
\begin{itemize}
    \item If $P_{j,h}(\lambda)$ is defined, then $\D{P,j,h}(\lambda)=d+1$, where $d$ is the unique integer satisfying  $\varrho_{j,h}(\lambda) = i - dh$.
    \item If $Q_{j,h}(\lambda)$ is defined, then $\D{Q,j,h}(\lambda)=d$, where $d$ is the unique integer satisfying $\vartheta_{j,h}(\lambda) = j + d(h-1)$.
\end{itemize}
We adopt the convention that \(\D{P,j,h}(\lambda)\) (resp. \(\D{Q,j,h}(\lambda)\)) is defined only when \(P_{j,h}(\lambda)\) (resp. \(Q_{j,h}(\lambda)\)) is.

\smallskip
\noindent
We refer to $\D{P,j,h}(\lambda)$ (resp. $\D{Q,j,h}(\lambda)$) as the degree of the operator $P_{j,h}$ (resp. $Q_{j,h}$) evaluated at $\lambda$.
\end{definition}


\begin{remark}\rm
Roughly speaking, $\D{P,j,h}(\lambda)$ measures the cost of passing from $\lambda$ to $P_{j,h}(\lambda)$ in the following sense. 
Recall that
$P_{j,h}(\lambda)=\lambda-\alpha_{\varrho,j}$,
where $\varrho =\varrho_{j,h}(\lambda)$. 
Thus, obtaining $P_{j,h}(\lambda)$ from $\lambda$ amounts to subtracting $\alpha_{\varrho,j}$. 
We can write 
\begin{equation} \label{descomponer raiz larga en suma de height h}
    \alpha_{\varrho ,j}=\sum_{k=0}^{d}\alpha_{i-kh,\,j-kh},
\end{equation}
with  $d$ as in \Cref{def: degree of P and Q}. 
Consequently, the degree corresponds to the number of summands in \cref{descomponer raiz larga en suma de height h}; equivalently, it quantifies how many roots of height $h$  must be subtracted from $\lambda $ to reach $P_{j,h}(\lambda)$.
Similar considerations hold for $\D{Q,j,h}$, with the role of $h$ replaced by $h-1$. 

Lastly, we have a notational remark. 
To simplify notation one might be tempted to write the degree as $D(P_{j,h}(\lambda))$ rather than $D(P_{j,h})(\lambda)$. 
We avoid this, because it would misleadingly suggest that we are assigning a degree to the weight $P_{j,h}(\lambda)$ itself, rather than to the ``passage''  $\lambda \to P_{j,h}(\lambda)$ encoded by $\D{P,j,h}(\lambda)$.
\end{remark}

We can extend the definition of degree for any composition of $P$ and $Q$ operators. 

\begin{definition} \rm
    Let $\lambda \in X$. 
    Let $Y=(Y_1,\ldots ,Y_{\ell})$ be a sequence of $P$ and $Q$ operators. 
    Suppose that $Y_{\ell} \cdots Y_1(\lambda)$ is defined.
    We set $\lambda_0=\lambda$ and $\lambda_k =Y_k\cdots Y_{1}(\lambda)$, for $1\leq k < \ell$.
    Then, we define the degree of the operator $Y$ evaluated at $\lambda$ is defined as
    \begin{equation}
        D(Y)(\lambda) = \sum_{k=1}^{\ell} D(Y_k)(\lambda_{k-1}). 
    \end{equation}
\end{definition}




We are now in position to state the decomposition mentioned at the beginning of this section.
 

\begin{proposition}\rm\label{prop: second version}
    Let $\lambda\in X^{+}$,  $\alpha_{i,j}\in \Phi^{\geq 2}$ and $h=j-i+1$ .  Let $k$ be an integer with $1\leq k <h$ satisfying  $\langle \lambda , \alpha_r \rangle  =0$, for all $j-k+1\leq r \leq j$. Then,  we have    
    \begin{equation}\label{eq: second version}
     \!  \bfM_{\lambda}^{\succeq \alpha_{i,j}}\! \!  = \! \bfM_{\lambda}^{\succ \alpha_{i,j}} \!  - q^{D(P^{k}_{j,h})(\lambda)+1} \bfM_{P^k_{j,h}(\lambda) - \alpha_{i-k,j-k}}^{\succ \alpha_{i,j}}\! \!  -  \displaystyle\sum_{r=1}^{k} q^{D(v^{r}_{j,h})(\lambda)} \! \left( \bfM_{v^{r}_{j,h}(\lambda)}^{\succ \alpha_{i,j}}\!  -\!  q\bfM_{v^{r}_{j,h}(\lambda)-\alpha_{i-r,j-r}}^{\succ \alpha_{i,j}}\right).
    \end{equation}
\end{proposition}

% \David{No s\'e si este convencion la hicimos antes}
\begin{remark}\rm\label{remark operator not defined}
       We emphasize that if any $P$-operator or $v^r$-operator is not defined, the corresponding $\bfM$-element is set  equal to zero. 
\end{remark}
    




%%%%%%%%%%%%%%%%%%%%%%%%%%%%%%%%%%%%%%%%%%%%%%%%%%%%%%%%%%%%%%%%%%%%%%%%%%%%%%%%%%%%%%%
\subsection{Identities involving $P$ and $Q$ operators}


In this section, building on the results of \Cref{section identities to prove first inverse decomposition}, we develop the tools required to prove \Cref{prop: second version}. 
That proof is presented at  the end of this section.


\begin{lemma}\rm \label{lem: free Q1} 
    Let $K \subset \Phi^{\geq 2}$, $\alpha_{i,j} \in \Phi^{\geq 2}$, and  $h = j - i + 1$. 
    Let $(r,p)$ be integers such that $i + 1 < r \leq p \leq j$  and $h+1 \leq r$. 
    Let $k\geq 0$ be  such that $p'\coloneqq p+k(h-1) \leq n$.
    Let  $Y = K \cup \{\alpha_{r-h, r-1}\}$ and $\lambda \in X^+_{p'+h-1,h}(R)$. 
\begin{enumerate}[(a)]
    \item   \label{item a} Suppose that $\OQ{p',h}(\lambda)$ is defined and let  $\vartheta \coloneqq \vartheta_{p',h}(\lambda)$ be the integer associated to $\OQ{p',h}(\lambda)$. 
    
    We define $d$ to be the unique integer satisfying $\vartheta = p +(d+1)(h-1)$.
    \item  \label{item b} Suppose that  $\OQ{p',h}(\lambda)$ is not defined.
    
    We define $d$ to be  the largest  integer such that $p + d(h - 1) \leq n$. 
\end{enumerate}
    In both cases we define
    $$T = \bigcup_{m = 0}^{d} [r + m(h - 1), p + m(h - 1)],$$
    and assume that $\pint{\lambda}{\alpha_t} = 0$ for all $t \in T \setminus \{p'\}$ and $\pint{\lambda}{\alpha_{p'}} = -1$. 
    Then, if $K \sim_{T} \Phi^{\succ \alpha_{i,j}}$, we have
    \begin{equation}
        \bfM_{\lambda}^{Y} = \left\{ \begin{array}{ll}
         q^{\D{Q,p',h}(\lambda)}\bfM_{Q_{p',h}(\lambda)}^{Y},    & \mbox{if } \OQ{p',h}(\lambda) \mbox{ is defined;}    \\
          0,    & \mbox{if } \OQ{p',h}(\lambda) \mbox{ is not defined.} 
        \end{array} \right.
    \end{equation}
\end{lemma}

\begin{proof}
We first address case \ref{item a}.  
That is, we assume that $\OQ{p',h}(\lambda)$ is defined. 
By definition, $\OQ{p',h}(\lambda) = \lambda - \alpha_{p'+1,\vartheta}$, and $\lambda - \alpha_{p'+1,\vartheta}\in X^{+}_{p',h}$
% is $R$-dominant at position $p'$ and height $h$. 
Moreover, by the definition of $\D{Q,p',h}(\lambda)$, we obtain 
\begin{equation}
    \D{Q,p',h}(\lambda) = \frac{\vartheta -p'}{h-1} = d - k + 1.
\end{equation}  
Applying \Cref{Claim: free Q} to $K$, $\alpha_{i,j}$, $(r,p,k)$ and $d$, it follows that 
\begin{equation}
    \bfM_{\lambda}^Y = q^{d-k+1} \bfM_{\lambda - \alpha_{p'+1,\vartheta}}^{Y} = q^{\D{Q,p',h}(\lambda)}\bfM_{Q_{p',h}(\lambda)}^{Y},
\end{equation}


We now treat case \ref{item b}.  
Thus, we assume that $\OQ{p',h}(\lambda)$ is not defined. 
By definition, $k \leq d$. If $k = d$, then the maximality of $d$ and the second case of \eqref{eq: M element with a -1 in p'} in \Cref{lem: right solving} imply that $\bfM_{\lambda}^{Y} = 0$.

Thus, we may assume that $k < d$.  
In this situation, we apply \Cref{Claim: free Q} to $K$, $\lambda$, $k$, and $d - 1$, yielding $\bfM_{\lambda}^{Y} = q^{d-k}\bfM_{\mu}^{Y}$, where $\mu = \lambda - \alpha_{p'+1,p'+(d-k)(h-1)}$ and $p'+(d-k)(h-1) = p+d(h-1)$ .  
Observe that via a simple calculation we have that  $\pint{\mu}{\alpha_{p+d(h-1)}} = -1$ and $\pint{\mu}{\alpha_{t}} = 0$ for $t \in T\setminus \{p + d(h-1)\}$.

By the maximality of $d$ and the second case of  \Cref{lem: right solving} applied to $K$, $\alpha_{i,j}$, $(r,p)$, $\mu$ and $T$, we obtain $\bfM_{\mu}^{Y} = 0$. 
Therefore, $\bfM_{\lambda}^Y = 0$, as required.
\end{proof}







\begin{lemma}\rm \label{lem: free Q2}
Let $K \subset \Phi^{\geq 2}$, $\alpha_{i,j} \in \Phi^{\geq 2}$ and let $h = j - i + 1$.
Suppose that $1 \leq i+1 - h$.
Let $p$ be an integer such that $i + 1 \leq p \leq j$. 
Let $k\geq 0$ be such that $p'\coloneqq p+k(h-1) \leq n$. 
Let $Y = K \cup \{\alpha_{i-h+1, i}\}$ and  $\lambda \in X_{p'+h-1,h}^+(R)$.

\begin{enumerate}[(a)]
    \item  \label{item a Q} Suppose that $\OQ{p',h}(\lambda)$ is defined. 
    Let $\vartheta\coloneqq \vartheta_{p',h}(\lambda)$ be the integer associated to $\OQ{p',h}(\lambda)$.  
    
    We define $d$ to be the unique integer satisfying $\vartheta = p+ (d+1)(h-1)$.
    
    \item  \label{item b Q} Suppose that $\OQ{p',h}(\lambda)$ is not defined.
    
    We define $d$ to be the largest integer such that $p + d(h - 1) \leq n$.
\end{enumerate}

In both cases we define
\begin{equation}
    T = \left( \bigcup_{m = 0}^{d-1} [i+1 + m(h - 1), j + m(h - 1)] \right) \cup [i+1+d(h-1),p+d(h-1)] 
\end{equation}
and assume that $\pint{\lambda}{\alpha_t} = 0$ for all $t \in T \setminus \{p'\}$, $\pint{\lambda}{\alpha_{p'}} = -1$.

Furthermore, if $k<d$ we also assume that $\pint{\lambda}{\alpha_{p'+1}} =1$.
If  $K \sim_{T} \Phi^{\succ \alpha_{i,j}}$  then we have
\begin{equation}
        \bfM_{\lambda}^{Y} = \left\{ \begin{array}{ll}
         q^{\D{Q,p',h}(\lambda)}\bfM_{Q_{p',h}(\lambda)}^{Y},    & \mbox{if } \OQ{p',h}(\lambda) \mbox{ is defined;}    \\
          0,    & \mbox{if } \OQ{p',h}(\lambda) \mbox{ is not defined.} 
        \end{array} \right.
\end{equation}
\end{lemma}

\begin{proof}
   The result follows by mimicking the proof of  \Cref{lem: free Q1},  using \Cref{Claim: free Q2} instead of \Cref{Claim: free Q}.
\end{proof}

\begin{lemma}\rm\label{lem: free P}  
Let $K \subset \Phi^{\geq 2}$, $\alpha_{i,j} \in \Phi^{\geq 2}$ and $h = j - i + 1$. 
Let $p$ be an integer such that $i -h +1< p \leq i$. 
Let $k\geq 0$ be such that $p'\coloneqq p-kh\geq 1$. 
Let $\lambda\in X_{p'-h,h}^+(L)$.
\begin{enumerate}[(a)]
    \item \label{item a P} Suppose that $P_{p'-1,h}(\lambda)$ is defined and let $\varrho= \varrho_{p'-1,h}(\lambda)$ be the integer associated to $P_{p'-1,h}(\lambda)$.
    
    We define $d$ to be the unique integer satisfying $\varrho = p-(d+1)h$.
 \item \label{item b P} Suppose that $P_{p'-1,h}(\lambda)$ is not defined.
    
     We define $d$ to be the largest integer such that $p-dh\geq 1$.
 \end{enumerate}
In both cases, we define $\displaystyle T=\bigcup_{m=0}^{d}[p-mh,i-mh]$
and assume that $\pint{\lambda}{\alpha_{t}}=0$ for all $t\in T\setminus\{p'\}$, $\pint{\lambda}{\alpha_{p'}}=-1$. 

If  $K \sim_{T} \Phi^{\succ \alpha_{i,j}}$ then we have
   \begin{equation}
        \bfM_{\lambda}^{K} = \left\{ 
        \begin{array}{ll}
         q^{\D{P,p'-1,h}(\lambda)}\bfM_{P_{p'-1,h}(\lambda)}^{K},    &  \mbox{if } P_{p'-1,h}(\lambda)  \mbox{ is defined;} \\
           0,  & \mbox{if } P_{p'-1,h}(\lambda)  \mbox{ is  not defined.}
        \end{array}
        \right.
   \end{equation}
\end{lemma}






\begin{proof}
We first address case \ref{item a P}.  
Assume that $\OP{p'-1,h}(\lambda)$ is defined. 
By definition, $\OP{p'-1,h}(\lambda) = \lambda - \alpha_{\varrho,p'-1}$, and $\lambda - \alpha_{\varrho,p'-1} \in X_{p'-h+1,h}(L)$
% is $L$-dominant at position $p'$ and height $h$. 
Moreover, by the definition of $d$, we deduce that $\D{P,p'-1,h}(\lambda) = d - k + 1$.  
Applying \Cref{Claim: free P}, it follows that 
\begin{equation}
    \bfM_{\lambda}^K = q^{d-k+1} \bfM_{\lambda - \alpha_{\varrho,p'-1}}^{K}.
\end{equation}
% \David{NO hay $Y$ aqu\'i!! No deber\'ia ser $K$. Revisar en toda la demo.}


\medskip
We now consider case \ref{item b P}.  
Assume that $\OP{p'-1,h}(\lambda)$ is not defined. 
By definition, $k \leq d$. 
If $k = d$ then the maximality of $d$ and the second case of \eqref{eq: left solving} in \Cref{lem: left solving} applied to $K$, $\alpha_{i,j}$, $p$, $\lambda$ and $d$ imply that $\bfM_{\lambda}^{K} = 0$.
Thus, we may assume that $k < d$.  
In this setting we apply \Cref{Claim: free P} to $K$, $\alpha_{i,j}$, $\lambda$, $k$, and $d - 1$, yielding $\bfM_{\lambda}^{K} = q^{d-k}\bfM_{\mu}^{K}$,
% \David{Creo que aca es el mismo error que en el caso (b) del lemma 4.17. Revisar!!!} 
where $\mu = \lambda - \alpha_{p'-(d-k)h,p'-1}$ and $p'-(d-k)h = p-dh$.  
Observe that via a simple calculation we have that $\pint{\mu}{\alpha_{p-dh}} = -1$ and $\pint{\mu}{\alpha_{t}} = 0$ for $t \in T\setminus \{p - d(h-1)\}$.
By maximality of $d$ and the second case of \Cref{lem: left solving} applied to $K$, $\alpha_{i,j}$, $p$ $\mu$, $d$ and $T$ it follows that
\begin{equation}
    \bfM_{\mu}^{K} = 0.
\end{equation}
Therefore, $\bfM_{\lambda}^K = 0$, as required.
\end{proof}






\begin{lemma}\rm\label{lem: pre completar v}
    Let $\alpha_{i,j} \in \Phi^{\geq 2}$ and $h=j-i+1$.
    Let $r$ and $x$ be integers such that $i+1 \leq r \leq j$, $h+1\leq r$ and $1\leq x \leq j-r+1$. 
    Let  $\lambda\in X^+_{r+1}(R)$ be such that $\pint{\lambda}{\alpha_{t}} =0$ for $t\in [r+1,j]$, $\pint{\lambda}{\alpha_{r}} =-1$ and $\pint{\lambda}{\alpha_{j+1}} \geq 1$. 
    If  $Y= \Phi^{\succ \alpha_{i,j}}\cup \{\alpha_{r-h,r-1}\}$ then
        \begin{equation}\label{eq: pre completar v}
            \bfM_{\lambda}^{Y} = \left\{\begin{array}{ll}
                q^{D(Q^{-x}_{r,h})(\lambda)}\bfM_{Q^{-x}_{r,h}(\lambda)}^{Y}, &\mbox{if } Q^{-x}_{r,h}(\lambda) \mbox{ is defined}; \\
                0, & \mbox{if } Q^{-x}_{r,h}(\lambda)\mbox{ is not defined}.
            \end{array}\right.
        \end{equation}
\end{lemma}



\begin{proof}
We only prove the statement when $Q^{-x}_{r,h}(\lambda)$ is defined, the other case being analogous. 

For $0\leq m \leq x$ we define $\mu_{m} = Q^{-m}_{r,h}(\lambda)$.
For $1\leq m \leq x$ we define $\vartheta_m=\vartheta_{r+m-1,h}(\mu_{m-1})$ and $d_m$ to be the unique integer satisfying $\vartheta_m =r+m-1+(d_m+1)(h-1) $.  

We split the proof in two cases. 

\begin{description}
    \item[Case A]  $r>i+1 $. 

    We proceed by induction on $x$. 
    If $x=1$ then $Q^{-x}_{r,h}(\lambda)=Q_{r,h}(\lambda)$. 
    Thus, the result reduces to \Cref{lem: free Q1} applied to 
    $K=\Phi^{\succ \alpha_{i,j}}$ and $p'=p=r$.
    Indeed, using the notation in that lemma we have 
    \begin{equation}
     T= \bigcup_{m=0}^{d_1}        \{ r+m(h-1)   \}.   
    \end{equation}
     Since $\lambda \in X^{+}_{r+1}(R)$,  the minimality of $\vartheta_1$ implies  $\pint{\lambda}{\alpha_{t}}=0$ for all $t\in T\setminus \{r\}$. 
     Thus we are under the hypothesis of \Cref{lem: free Q1}, which allows us to conclude.

We now assume that $x\geq 2$ and suppose that \eqref{eq: pre completar v} holds for $x-1$. 
This is, 
   \begin{equation}\label{eq: hp completar v hasta p}
     \displaystyle   \bfM_{\lambda }^{Y} = q^{D(Q^{-(x-1)}_{r,h})(\lambda)}\bfM_{\mu_{x-1}}^{Y}.
    \end{equation}
Let 
 \begin{equation}
        T  = \bigcup_{m=0}^{d_x}[r+m(h-1),r+x-1+m(h-1)].
\end{equation}

By definition,  we have  $\displaystyle \mu_{x-1} = \lambda - \sum_{m=1}^{x-1}\alpha_{r+m,\vartheta_{m}}$.
A direct computation, using \Cref{remark producto interno}, together with \Cref{lem: ceros por Qs} applied to $\mu_{x}$, 
imply that   $\pint{\mu_{x-1}}{\alpha_{r+x-1}} = -1$ and $\pint{\mu_{x-1}}{\alpha_{t}} = 0$, for all $t\in T\setminus\{r+x-1\}$.
Then, we can apply \Cref{lem: free Q1} to $K=\Phi^{\succ \alpha_{i,j}}$, $r$, $p=r+x-1=p'$ and $\mu_{x-1}$ to get
    \begin{equation}\label{eq: completar v hasta p}
        \bfM_{\mu_{x-1}}^{Y} = q^{D(Q_{r+x-1,h})(\mu_{x-1})}\bfM_{Q_{r+x-1,h}(\mu_{x-1})}^{Y}  = q^{D(Q_{r+x-1,h})(\mu_{x-1})}\bfM_{\mu_x}^{Y}.
    \end{equation}
    We conclude by combining \eqref{eq: hp completar v hasta p}, \eqref{eq: completar v hasta p} and  $D(Q^{-x}_{r,h})(\lambda) = D(Q^{-(x-1)}_{r,h})(\lambda) + D(Q_{r+x-1,h})(\mu_{x-1})$.



\item[Case B]  $r=i+1 $. 

Using the assumption $\pint{\lambda}{\alpha_{j+1}}\ge 1$ and an inductive argument (similar to the one in the proof of \Cref{lem: thetas para el caso h-1}), we obtain $\vartheta_m=j+m$ for all $1\le m\le x$.
With this in hand, the result follows by the same inductive method as in the proof of \textbf{Case A}, replacing \Cref{lem: free Q1} with \Cref{Claim: free Q2}.
\end{description}
\end{proof}





\begin{corollary}\rm \label{lem: v existe o no existe}
    Let $\lambda\in X^{+}$,  $\alpha_{i,j}\in \Phi^{\geq 2}$ and $h=j-i+1$ .  
    Let $s$ be an integer with $1 \leq s \leq h-1$ satisfying  $\langle \lambda , \alpha_t \rangle  = 0$, for all $j-s+1\leq t \leq j$. 
    Set $\lambda_{s} = P^s_{j,h}(\lambda)$.
   If $\lambda_{s}$ is defined then we have
    \begin{equation}\label{eq: v existe o no existe}
        \hspace{-10pt}\bfM_{\lambda_{s}}^{\succ \alpha_{i,j}} = \left\{\begin{array}{ll}
        q\bfM_{\lambda_{s}-\alpha_{i-s,j-s}}^{\succ \alpha_{i,j}} + q^{D(Q^{-s}_{j-s+1,h})(\lambda_{s})}\left(\bfM_{v^s_{j,h}(\lambda)}^{\succ \alpha_{i,j}} - q\bfM_{v^s_{j,h}(\lambda)-\alpha_{i-s,j-s}}^{\succ \alpha_{i,j}}\right), & \mbox{if } v^s_{j,h}(\lambda) \mbox{ is defined};\\
        q\bfM_{\lambda_{s}-\alpha_{i-s,j-s}}^{\succ\alpha_{i,j}},  & \mbox{otherwise}.
        \end{array}\right.
    \end{equation}    
\end{corollary}


\begin{proof}
By definition of $\lambda_s$ and the hypothesis on $\lambda$  we have   $\pint{\lambda_{s}}{\alpha_{j-s+1}} = -1$ and $\pint{\lambda_{s}}{\alpha_{t}} = 0$ for all $t\in [j-s+2,j]$. 
Furthermore, since $\lambda \in X^+$ it follows that $\lambda_{s}\in X^+_{j-s+2}(R)$. 


We recall that
\begin{equation}
  v^{s}_{j,h}(\lambda) = Q^s_{j,h}P^{s}_{j,h}(\lambda) = Q^s_{j,h}(\lambda_s)= Q^{-s}_{j-s+1,h}(\lambda_s ).   
\end{equation}

In particular,  $Q^{-s}_{j-s+1,h}(\lambda_s )$ is defined if and only if   $v^{s}_{j,h}(\lambda)$ is. 

Let  $Y = \Phi^{\succ \alpha_{i,j}}   \cup \{ \alpha_{i-s,j-s} \}$.
By applying \Cref{lem: pre completar v} to $\alpha_{i,j}$, $(r,p) = (j-s+1,s)$ and $\lambda_{s}$, it follows that
\begin{equation}\label{eq: 4.21 pre completar v}
            \bfM_{\lambda_s}^{Y} = \left\{\begin{array}{ll}
               q^{D(Q^{-s}_{j-s+1,h})(\lambda_{s})}\bfM_{v^s_{j,h}(\lambda )}^{Y}, &\mbox{if } v^s_{j,h}(\lambda) \mbox{ is defined}; \\
                0, & \mbox{if } v^s_{j,h}(\lambda) \mbox{ is not  defined}.
            \end{array}\right.
        \end{equation}
Then, the corollary follows by expanding according to the rule $\bfM_{\mu}^{Y} = \bfM_{\mu}^{\succ \alpha_{i,j}} - q \bfM_{\mu -\alpha_{i-s,j-s}}^{\succ \alpha_{i,j}}$
\end{proof}



\begin{lemma}\rm \label{lem: los 3 casos de second version}
    Let $\lambda\in X^{+}$,  $\alpha_{i,j}\in \Phi^{\geq 2}$ and $h=j-i+1$ .  
    Let $s$ be an integer with $0 \leq s < h-1$ satisfying  $\langle \lambda , \alpha_t \rangle  = 0$, for all $j-s\leq t \leq j$. 
    Let $\lambda_{s} = P^{s}_{j,h}(\lambda)$ and $\lambda_{s+1}=P^{s+1}_{j,h}(\lambda)$.
    We also define 
    \begin{equation}
       u=\D{P,j-s,h}(\lambda_{s})+1 \mbox{ and } v = D(Q^{-s}_{j-s+1,h})(\lambda_s) + \D{P,j-s,h}(\lambda_{s}).\footnote{If the right-hand side is defined. } 
    \end{equation}
    If $\lambda_s$ is defined then 
\begin{equation}  \label{eq: casicasi First}
    q\bfM_{\lambda_{s} - \alpha_{i-s,j-s}}^{\succ \alpha_{i,j}} =   \begin{cases}
      q^{u}\bfM_{\lambda_{s+1}-\alpha_{i-(s+1),j-(s+1)}}^{\succ \alpha_{i,j}} + q^{v}\left(\bfM_{v^{s+1}_{j,h}(\lambda)}^{\succ \alpha_{i,j}} - q\bfM_{v^{s+1}_{j,h}(\lambda)-\alpha_{i-(s+1),j-(s+1)}}^{\succ \alpha_{i,j}}\right)  &  \\
     q^{u}\bfM_{\lambda_{s+1}-\alpha_{i-(s+1),j-(s+1)}}^{\succ\alpha_{i,j}}   &   \\
     0, &  
    \end{cases}
\end{equation}
where the three cases correspond to:  $v^{s+1}_{j,h}(\lambda)$ is defined,  $\lambda_{s+1}$ is defined but $v^{s+1}_{j,h}(\lambda)$ is not, and $\lambda_{s+1}$  is not defined, respectively.
\end{lemma}

\begin{remark}\rm
    Since $v^{s+1}_{j,h}(\lambda)= Q^{s+1}_{j,h}P^{s+1}_{j,h}(\lambda) =Q^{s+1}_{j,h}(\lambda_{s+1})$, it follows that if $\lambda_{s+1}$ is not defined then $v^{s+1}_{j,h}(\lambda )$ is not defined. 
\end{remark}

    


\begin{proof}
    Suppose $\lambda_{s+1} = P_{j-s,h}(\lambda_{s})$ is defined. 
    We begin by noticing that if $\pint{\lambda_{s}}{\alpha_{i-s}} \geq 1$, then by definition of $P$-operator, it follows that $\lambda_{s+1}= \lambda_{s} - \alpha_{i-s,j-s}$.
    Therefore, $\bfM_{\lambda_{s} - \alpha_{i-s,j-s}}^{\succ \alpha_{i,j}} = \bfM_{\lambda_{s+1}}^{\succ \alpha_{i,j}}$ and the result follows by \Cref{lem: v existe o no existe}. 
    
    Thus, we assume that $\pint{\lambda_{s}}{\alpha_{i-s}} = 0$.
    By \Cref{lem: ceros por Ps} applied to $\lambda\in X^{+}$, $\alpha_{i,j}$ and $k=s+1<h$ we get
    \begin{equation}
        \pint{\lambda_{s}}{\alpha_{t}} = 0 \quad \mbox{for all}\; t\in T=\bigcup_{m=0}^{d} [i-s-mh,i-mh],
    \end{equation}
    where $d$ is the unique integer satisfying $\varrho_{j-s,h}(\lambda_{s}) = i-s-(d+1)h$. 
    It follows from \Cref{remark producto interno} that
    \begin{equation}  \label{eq: lemma casicasi 0 y-1}
        \pint{\lambda_{s} - \alpha_{i-s,j-s}}{\alpha_{t}}  = \begin{cases}
            0, & \mbox{if } t \in T \setminus \{i-s\};\\
           -1, & \mbox{if } t=i-s.
        \end{cases}
    \end{equation}
    Hence, we can apply \Cref{lem: free P} to $K=\Phi^{\succ \alpha_{i,j}}$, $\alpha_{i,j}$, $p=p'=i-s$, $T$ and $\lambda_{s}-\alpha_{i-s,j-s}$, to get
    \begin{equation}\label{eq: p existe en second version 1}
        q\bfM_{\lambda_{s} - \alpha_{i-s,j-s}}^{\succ \alpha_{i,j}} = q^{\D{P,i-s-1,h}(\lambda_{s} - \alpha_{i-s,j-s})+1}\bfM_{P_{i-s-1,h}(\lambda_{s}-\alpha_{i-s,j-s})}^{\succ \alpha_{i,j}}
        = q^{\D{P,j-s,h}(\lambda_{s})}\bfM_{\lambda_{s+1}}^{\succ \alpha_{i,j}},
    \end{equation}
    where the second equality holds since 
    \begin{equation}
        P_{i-s-1,h}(\lambda_{s}-\alpha_{i-s,j-s}) = P_{j-s,h}(\lambda_{s})\quad \mbox{ and } \quad \D{P,i-s-1,h}(\lambda_{s} - \alpha_{i-s,j-s})+1 = \D{P,j-s,h}(\lambda_{s}).
    \end{equation}
   Thus, using  \Cref{lem: v existe o no existe} for $s+1$ to expand the right-hand side of  \eqref{eq: p existe en second version 1} we obtain the first two cases of \cref{eq: casicasi First}. 

   We now  suppose that $\lambda_{s+1}$ is not defined.  
    By \Cref{lem: ceros por Ps} applied to $\lambda$, $\alpha_{i,j}$ and $k=s$ it follows that $\pint{\lambda_{s-1}}{\alpha_{t}} = 0$ for all $t\in T$, where
    \begin{equation}
        T=\bigcup_{m=0}^{d}[i-s+1-mh,i-mh],
    \end{equation}
    and $d$ is the unique integer satisfying $\varrho_{s} = \varrho_{j-s+1,h}(\lambda_{s-1}) = i-s+1-(d+1)h$.
    On the other hand, since $\lambda_{s+1}$ is not defined, it follows by \Cref{lem: desigualdad de varrhos} that $\varrho_{s} = 1$.
    In other words, $\lambda_{s} = \lambda_{s-1} -\alpha_{1,j-s+1}$.
   By \Cref{remark producto interno} we conclude that $\pint{\lambda_{s}}{\alpha_{t}}=0$ for all $t\in T$.
    Moreover, given that $\lambda_{s+1}$ is not defined, by the definition of the  $P$ operator we have  $\pint{\lambda_{s}}{\alpha_{t}} =0$ for all $t\in \{i-s-mh\mid 0\leq m\leq d\}$.
    Thus, we conclude that $\pint{\lambda_{s}}{\alpha_{t}} =0$ for all $t\in T'$, where
    \begin{equation}
        T'=\bigcup_{m=0}^{d} [i-s-mh,i-mh].
    \end{equation} 
 Using \Cref{remark producto interno} we obtain \eqref{eq: lemma casicasi 0 y-1} for all $t\in T'$. 
 As before, we can apply \Cref{lem: free P} to $K=\Phi^{\succ \alpha_{i,j}}$, $\alpha_{i,j}$, $p=p'=i-s$, $T'$ and $\lambda_{s}-\alpha_{i-s,j-s}$, to get $\bfM_{\lambda_{s}-\alpha_{i-s,j-s}}^{\succ \alpha_{i,j}}=0$, thus proving the third case in \Cref{eq: casicasi First}. 
\end{proof}

We are now in position to prove the First Inverse Decomposition given in \Cref{prop: second version}. 


\begin{myproof}[Proof of \Cref{prop: second version}]
We begin by emphasizing that in \cref{eq: second version} the $\bfM$-elements indexed by undefined weights are set equal to $0$ by convention.
    Let $1\leq k< h$ satisfy the hypothesis of the \Cref{prop: second version}, i.e. $\pint{\lambda}{\alpha_r}=0$ for all $j-k+1\leq r \leq j$.  
    We  proceed  by  induction on $k$.
    Consider first the case $k = 1$.
   % Notice that in this case we must have $r=j$. 
    By the definition of $\bfM$-elements, we have
    \begin{equation}\label{eq: second version -1}
    \bfM_{\lambda}^{\succeq \alpha_{i,j}} = \bfM_{\lambda}^{\succ \alpha_{i,j}} - q\bfM_{\lambda-\alpha_{i,j}}^{\succ \alpha_{i,j}}.
    \end{equation}
  By applying  \Cref{lem: los 3 casos de second version} for $s=0$ to rewrite $q\bfM_{\lambda-\alpha_{i,j}}^{\succ \alpha_{i,j}}$ we obtain \eqref{eq: second version} for $k=1$,  as we wanted.
 
We now assume that \cref{eq: second version} holds for some $1\leq k < h-1$.
Let $\lambda_{s} = P^s_{j,h}(\lambda)$ for $1\leq s\leq k+1$. 
We also let $\mathcal{RH}_k$ and $\mathcal{RH}_{k+1}$ be the right-hand side of \eqref{eq: second version} for $k$ and $k+1$, respectively. 
We must show that $\bfM_{\lambda}^{\succeq \alpha_{i,j}}=\mathcal{RH}_{k+1}$. 
By our inductive hypothesis, this amounts to show that $\mathcal{RH}_k=\mathcal{RH}_{k+1}$. 
This turns out to be  equivalent to 


\begin{equation} \label{eq: turns out to be equivalent}
  q \bfM_{\lambda_{k} - \alpha_{i-k,j-k}}^{\succ \alpha_{i,j}} = q^{u} \bfM_{\lambda_{k+1} - \alpha_{i-(k+1),j-(k+1)}}^{\succ \alpha_{i,j}} 
  + q^{ v   }  
  \left( \bfM_{v^{k+1}_{j,h}(\lambda)}^{\succ \alpha_{i,j}} - q \bfM_{v^{k+1}_{j,h}(\lambda) - \alpha_{i-(k+1),j-(k+1)}}^{\succ \alpha_{i,j}} \right),
\end{equation}
where $u=D(P_{j-k})(\lambda_k) + 1$ and $v=D(Q^{k+1}_{j,h})(\lambda_{k+1}) +D(P_{j-k,h} )(\lambda_k)  $.

Due to the conventions mentioned at the beginning of this proof,  \cref{eq: turns out to be equivalent} follows by \Cref{lem: los 3 casos de second version}. 
This completes the induction and the proof of the First Inverse Decomposition.
\end{myproof}


